% =========================================================
% VII. DISCUSSION AND LIMITATIONS
% =========================================================
\section{Thảo luận và hạn chế}
\label{sec:discussion_limitations}

Thực nghiệm có kiểm soát cho thấy truy xuất event--rule có cấu trúc và selective
verification cải thiện exact decision label. Full LegalSCM đạt 96,80\% so với
90,80\% của Causal Path RAG, với mức tăng theo cặp có ý nghĩa thống kê. Do hai
hệ thống có retrieval/path giống nhau và cả 17 Full-only gain đều xuất hiện sau
verifier commit, chênh lệch được khoanh vùng ở factual validation đã commit và
grounded direct-edge checking.

Kết quả không có nghĩa mọi khía cạnh đều tốt hơn. Answer F1 và ROUGE-L của Full
thấp hơn nhẹ; khác biệt theo cặp của hai chỉ số này và Citation F1 đều không có
ý nghĩa. Structured-label correctness và lexical answer overlap do đó là hai
mục tiêu khác nhau. Ba mẫu lỗi còn có prose mâu thuẫn với structured decision,
cho thấy cần constraint consistency riêng ở tầng sinh.

Kết quả closed-book LLM phụ thuộc mạnh vào prompt và task. Accuracy 0\% xuất
hiện vì cả 250 output thành công đều abstain bằng \textsc{uncertain}, trong khi
gold không có abstention class. Đây không phải bằng chứng Qwen3-8B không có tri
thức pháp luật. Phân bố gold 230/20 cũng khiến raw accuracy dễ bị thổi phồng;
vì vậy cần báo cáo kèm Balanced Accuracy, Macro-F1, class coverage và confusion
matrix.

Quan trọng nhất, thực nghiệm này không xác nhận structural mediator
intervention. Hai mươi mẫu âm là direct-edge counterexample và Full run đã lưu
thực thi 0 phép $do(X_M=0)$. Điểm subset 95\% đo grounded direct-edge topology.
Hard intervention đã được cài đặt và phân biệt chính thức với phép xóa nút
$G^{-M}$, nhưng cần benchmark mới với câu hỏi removal/necessity tường minh trước
khi kết luận về hiệu quả. Tương tự, engine hỗ trợ conjunctive body và alternative
disjunctive mechanisms, còn các quy tắc import hiện tại đều là unary positive;
hành vi AND/OR phong phú chưa được kiểm thử thực nghiệm.

Top-1 Exact Path chỉ đạt 38,56\%, so với Oracle Exact Path 75,42\%. Bảy trong
tám lỗi Full xảy ra sau missing-primary-path fallback, xác nhận first-hop
retrieval và path ranking vẫn là nút thắt. Ngoài ra, benchmark còn nhỏ, được
sinh từ đồ thị và dùng cùng nguồn cấu trúc với hệ thống. Cần rà soát chuyên gia,
câu hỏi pháp lý ngoài graph và đánh giá miền ngoài để bảo đảm external validity.

Cuối cùng, LegalSCM phụ thuộc vào rule extraction, event normalization,
endpoint/mediator grounding và graph construction. Đây là mô hình quy tắc quy
phạm có tính xác định, không phải causal discovery hoặc effect estimation. Toán
tử $do(\cdot)$ chỉ có nghĩa tương đối với các cơ chế đã mã hóa và không có tầng
abduction biến ngoại sinh của phân tích counterfactual cấp cá thể đầy đủ.

% =========================================================
% VIII. CONCLUSION
% =========================================================
\section{Kết luận}
\label{sec:conclusion}

LegalSCM-RAG kết hợp causal-path retrieval, suy luận quy tắc ba giá trị và
selective commit gate cho hỏi đáp pháp luật đa bước tiếng Việt. Trên benchmark
250 câu hiện tại, hệ thống tăng Decision Accuracy 6,0 điểm so với Causal Path
RAG với bằng chứng thống kê theo cặp, nhưng không cải thiện có ý nghĩa các chỉ
số answer hoặc citation.

Kết luận mạnh nhất được dữ liệu hỗ trợ là selective factual/direct-claim
verification cải thiện độ tin cậy quyết định trong môi trường có kiểm soát.
Nghiên cứu tiếp theo sẽ xây benchmark explicit mediator intervention và quy tắc
nhiều antecedent, cải thiện first-hop retrieval cùng endpoint grounding, ràng
buộc consistency answer--decision và tổ chức đánh giá chuyên gia trên câu hỏi
không sinh từ source graph.
